\documentclass[11pt,a4paper]{article}
\usepackage{fontspec}
\setmainfont{DejaVu Sans}[Scale=0.90]
\usepackage{amsmath,amssymb}
\usepackage[margin=2.2cm]{geometry}
\usepackage{booktabs,array,tabularx}
\usepackage{graphicx}
\usepackage{enumitem}
\usepackage{titlesec}
\usepackage{parskip}
\titlelabel{\thetitle.\quad}
\titleformat{\section}{\normalfont\large\bfseries}{\thesection.}{0.6em}{}
\titleformat{\subsection}{\normalfont\normalsize\bfseries}{\thesubsection}{0.6em}{}
\titlespacing*{\section}{0pt}{14pt}{5pt}
\titlespacing*{\subsection}{0pt}{10pt}{4pt}
\setlist[itemize]{leftmargin=1.4em,itemsep=1pt,topsep=3pt}

\begin{document}

\begin{center}
{\Large\bfseries Régimes de transport phononique\\et position de la condition de cécité}\\[5pt]
{\normalsize Note de partie III --- point de contrôle}\\[2pt]
{\small 13 septembre 2026}
\end{center}

\vspace{4pt}

\section{Objet}

Le transport phononique dans un film mince relève de trois régimes selon la hiérarchie des temps de
collision. La partie~IV établit par ailleurs une condition sous laquelle un milieu de
Guyer--Krumhansl devient thermiquement indiscernable d'un milieu de Fourier. La présente note situe
cette condition sur la carte des régimes.

\section{Variables réduites}

Le critère de Guyer pour la fenêtre hydrodynamique s'écrit
\[ \tau_N \;\ll\; \tau_B \;\ll\; \tau_R , \]
où $\tau_N$ est le temps des processus normaux, $\tau_R$ celui des processus résistifs, et $\tau_B$
celui associé à la diffusion aux frontières. Deux nombres sans dimension en découlent :
\[ x = \frac{\tau_R}{\tau_N} , \qquad\qquad y = \frac{\tau_B}{\tau_N} . \]

Le second admet une lecture directe. Avec $\tau_B = d/v$ et $\tau_N = \Lambda_N / v$, il vient
$y = d/\Lambda_N$ : l'épaisseur du film exprimée en libres parcours moyens des processus normaux.

Le critère devient alors $1 \ll y \ll x$, et la carte se partage en trois régions.

\begin{center}
\begin{tabular}{@{}ll@{}}
\toprule
\textbf{Condition} & \textbf{Régime} \\
\midrule
$y < 1$ & balistique : les frontières l'emportent sur les processus normaux \\
$y > x$ & diffusif : les processus résistifs l'emportent sur les frontières \\
$1 < y < x$ & hydrodynamique \\
\bottomrule
\end{tabular}
\end{center}

La fenêtre hydrodynamique est un coin dont le sommet se trouve en $(1,\,1)$. Sa largeur, en décades,
vaut $\log_{10} x$.

\section{Position de la condition de cécité}

La partie~IV établit qu'un milieu de Guyer--Krumhansl dont les deux temps coïncident,
$\tau_R = \tau_\ell$, produit exactement la réponse d'un milieu de Fourier, quelle que soit leur
valeur commune. Avec $\tau_\ell = \tfrac{9}{5}\tau_N$, cette condition s'écrit
\[ \tau_R = 1{,}8\,\tau_N , \qquad\text{soit}\qquad x = 1{,}8 . \]

C'est une droite verticale, et elle passe \textbf{juste à droite du sommet du coin}.

\begin{center}
\includegraphics[width=0.78\textwidth]{05_regime_map.png}
\end{center}

\begin{center}
\begin{tabular}{@{}cc@{}}
\toprule
$x = \tau_R/\tau_N$ & Largeur de la fenêtre \\
\midrule
$1{,}2$ & $0{,}08$ décade \\
$\mathbf{1{,}8}$ & $\mathbf{0{,}26}$ \textbf{décade} \\
$5$ & $0{,}70$ décade \\
$10$ & $1{,}00$ décade \\
$100$ & $2{,}00$ décades \\
\bottomrule
\end{tabular}
\end{center}

\subsection{Énoncé}

À l'abscisse de la cécité, la fenêtre hydrodynamique mesure un quart de décade. Elle n'en atteint une
qu'à $x = 10$.

Un matériau qui refroidit voit $\tau_R/\tau_N$ croître, les processus umklapp gelant plus rapidement
que les processus normaux. Il traverse donc la droite de cécité \textbf{au moment précis où la
fenêtre hydrodynamique s'ouvre}.

\subsection{Conséquence}

Au voisinage de l'entrée dans le régime hydrodynamique, un matériau devrait paraître thermiquement de
Fourier alors même qu'il cesse de l'être. La signature non-Fourier ne deviendrait détectable qu'une
fois la fenêtre franchement ouverte, soit à $x$ nettement supérieur à l'unité.

Cette conséquence est falsifiable : elle prédit, sur une mesure en température, une plage où l'écart
à Fourier passe par un minimum au lieu de croître de façon monotone.

\section{Position de l'AlN}

\subsection{Temps de relaxation}

Les temps de relaxation de l'AlN wurtzite à 300~K sont relevés sur la figure~7 de Ma, Li et Luo,
entre 2 et 10~THz. La lecture est faite à l'{\oe}il sur une échelle logarithmique et porte une
incertitude d'un facteur deux. Elle se valide toutefois elle-même : les exposants ajustés sur les
valeurs relevées donnent $-2{,}00$ pour $\tau_N$ de la branche LA et $-1{,}00$ pour celle des
branches TA, exactement les lois énoncées dans l'article.

\begin{center}
\begin{tabular}{@{}lcccc@{}}
\toprule
Branche & $\omega$ & $\tau_U$ & $\tau_N$ & $x = \tau_U/\tau_N$ \\
\midrule
LA & 2 THz & $7\times10^{5}$ ps & $1{,}0\times10^{4}$ ps & 70 \\
LA & 10 THz & $1{,}0\times10^{4}$ ps & $4\times10^{2}$ ps & 25 \\
TA & 2 THz & $4\times10^{5}$ ps & $2{,}0\times10^{3}$ ps & 200 \\
TA & 10 THz & $4\times10^{3}$ ps & $4\times10^{2}$ ps & 10 \\
\bottomrule
\end{tabular}
\end{center}

\textbf{L'abscisse de l'AlN à 300~K vaut donc $x = 10$ à $200$.} La droite de cécité étant à $1{,}8$,
le facteur de sécurité est de six au pire : \textbf{l'AlN n'est pas près de la cécité à température
ambiante}.

La plage d'un facteur vingt sur une seule décade de fréquence confirme quantitativement l'énoncé de
la section~5.2 : un matériau occupe une plage, non un point.

\subsection{Ordonnée et régimes}

Le libre parcours moyen des processus normaux vaut $\Lambda_N = v\,\tau_N$, soit $2{,}4$ à
$60\ \mu$m avec $v = 6000$~m/s. Cette longueur est très grande devant l'épaisseur d'un film mince.

\begin{center}
\begin{tabular}{@{}lcl@{}}
\toprule
Épaisseur & $y = \tau_B/\tau_N$ & Régime \\
\midrule
\textbf{500 nm} & $0{,}008$ à $0{,}21$ & \textbf{balistique} \\
5 $\mu$m & $0{,}08$ à $2{,}1$ & mixte \\
20 $\mu$m & $0{,}33$ à $8{,}3$ & mixte \\
500 $\mu$m & $8{,}3$ à $208$ & hydrodynamique \\
\bottomrule
\end{tabular}
\end{center}

\subsection{Conséquence}

\textbf{À 300~K, un film d'AlN submicronique est balistique.} Ni la loi de Fourier, ni celle de
Cattaneo, ni celle de Guyer--Krumhansl ne s'y appliquent : toutes trois sont des lois de milieu
continu, et l'épaisseur est trop faible pour qu'une description par équation de diffusion conserve un
sens.

Ce placement est cohérent avec la transition balistique vers diffusif mesurée par Hoque \emph{et al.}
sur des films d'AlN de 1,6 à 2440~nm.

Il en résulte que \textbf{l'épaisseur des films est la première caractéristique à établir}, avant la
bande de fréquences de la mesure. Sur un dépôt submicronique, l'extraction de propriétés par
inversion d'un modèle de diffusion repose sur un modèle hors de son domaine de validité.

\section{Portée et limites}

\subsection{La carte reste en variables réduites}

Les bandes portées sur la figure correspondent à une seule température, 300~K. La trajectoire de
l'AlN lorsque la température varie exige les coefficients d'un modèle de Callaway ajusté, ou le
modèle de Debye--Callaway modifié de Morelli qui les exprime par les paramètres de Grüneisen. Ni
l'un ni l'autre n'est établi à ce jour.

Les valeurs employées proviennent d'une lecture graphique et portent un facteur deux. Cela n'autorise
pas de conclusion au pourcent près ; cela suffit au placement des régimes, qui se joue sur des
décades.

\subsection{Le rapport des temps dépend de la fréquence}

Les calculs de premiers principes sur l'AlN wurtzite donnent, à basse fréquence,
$1/\tau_U \propto \omega^3$ pour les branches acoustiques, tandis que $1/\tau_N \propto \omega$ pour
les modes transverses et $\omega^2$ pour les longitudinaux.

Il en résulte $\tau_R/\tau_N \propto \omega^{-2}$ pour les modes transverses. \textbf{L'abscisse de
la carte n'est donc pas une constante du matériau} : les modes de basse fréquence se placent loin à
droite, ceux de haute fréquence à gauche. Un matériau occupe une plage, non un point.

La condition de cécité portant sur les temps moyennés de l'équation de Guyer--Krumhansl, sa
confrontation à un matériau réel exige de savoir comment cette moyenne est pondérée. Ce point n'est
pas tranché.

\subsection{Erreur propre du cadre}

Le modèle de Callaway s'écarte de la solution exacte de l'équation de Boltzmann de $+1\,\%$ dans le
plan et de $-12\,\%$ hors plan pour l'AlN. La direction hors plan étant celle de la configuration
étudiée, ce cadre porte une erreur propre de l'ordre de dix pour cent, antérieure à toute mesure.

\subsection{Ce qui est supposé}

Le critère de Guyer emploie un temps de frontière unique, $\tau_B = d/v$, sans facteur de forme ni
spécularité. Une frontière partiellement spéculaire allonge ce temps et déplace la limite balistique
vers le bas.

Les trois régions sont séparées par des inégalités strictes traitées ici comme des égalités. Les
frontières de la carte sont donc des transitions graduelles, non des lignes.

\section{Suite}

\begin{itemize}
\item Récupérer les coefficients $B_N$ et $B_U$ de l'AlN. Deux références à ouvrir :
Phys. Rev. B \textbf{90}, 035203, source des lois d'échelle ci-dessus ; et
Phys. Rev. Materials \textbf{4}, 044602, mesures sur monocristal de 130 à 480~K.
\item Tracer la trajectoire en température, pour plusieurs épaisseurs. La fonction correspondante
est écrite et attend ces coefficients.
\item Confronter aux mesures de transition balistique vers diffusif publiées sur films d'AlN de
1,6 à 2440~nm.
\end{itemize}

\vspace{4pt}
\noindent\small
Figure et calculs : \texttt{notebooks/05\_regime\_map.py}. Données matériau et limites :
\texttt{notes/conventions.md}, section~11. Condition de cécité : note de partie~IV, section~6.
Références à vérifier à la source avant citation.

\end{document}
